%% Đề 16 -- de-kiem-tra-cuoi-hoc-ki-1-toan-9-nam-2023-2024-so-gddt-lang-son.pdf
\deheader{ĐỀ SỐ 16}%
  {Đề kiểm tra Học kỳ — Năm học 2023--2024}%
  {Sở GD\&ĐT Lạng Sơn | Chỉ phần Đại số}
\addcontentsline{toc}{section}{Đề 16 — Đề kiểm tra Học kỳ — Năm học 2023--2024}

\begin{de_cau}{Câu 1}{2,0 điểm}
\begin{enumerate}[label=\alph*)]
\item Tính: $3\sqrt{8}-\sqrt{50}+\sqrt{18}$.
\item Cho $P=\dfrac{\sqrt{x}+1}{\sqrt{x}-1}-\dfrac{\sqrt{x}-1}{\sqrt{x}+1}-\dfrac{4\sqrt{x}-2}{x-1}$, $x\geq 0, x\neq 1$.
Rút gọn $P$ và tính $P$ khi $x=\frac{1}{4}$.
\end{enumerate}
\end{de_cau}

\begin{loi_giai_cau}
\textbf{a)} $6\sqrt{2}-5\sqrt{2}+3\sqrt{2}=4\sqrt{2}$.

\textbf{b)}
\[P=\frac{(\sqrt{x}+1)^2-(\sqrt{x}-1)^2-(4\sqrt{x}-2)}{(\sqrt{x}-1)(\sqrt{x}+1)}
=\frac{4\sqrt{x}-4\sqrt{x}+2}{x-1}=\frac{2}{x-1}.\]

Khi $x=\frac{1}{4}$: $P=\dfrac{2}{\frac{1}{4}-1}=\dfrac{2}{-\frac{3}{4}}=-\dfrac{8}{3}$.
\end{loi_giai_cau}

\begin{de_cau}{Câu 2}{2,0 điểm}
\begin{enumerate}[label=\alph*)]
\item Giải hệ phương trình: 
\[ \begin{cases}x+2y=4\\3x-y=5\end{cases} \].
\item Vẽ đồ thị hàm số $y=-x^2$ trên hệ trục tọa độ Oxy.
\end{enumerate}
\end{de_cau}

\begin{loi_giai_cau}
\textbf{a)} Từ (1): $x=4-2y$. Thế vào (2): $3(4-2y)-y=5\Rightarrow 12-7y=5\Rightarrow y=1, x=2$.

Nghiệm $(x;y)=(2;1)$.

\textbf{b)} Bảng giá trị $y=-x^2$:
\begin{center}
\begin{tabular}{c|ccccc}$x$&$-2$&$-1$&$0$&$1$&$2$\\\hline$y$&$-4$&$-1$&$0$&$-1$&$-4$\end{tabular}
\end{center}
Parabol có đỉnh $O(0;0)$, bề lõm quay xuống.
\end{loi_giai_cau}

\begin{de_cau}{Câu 3}{2,0 điểm}
Hai công nhân cùng làm một công việc, dự định hoàn thành trong 12 ngày. Sau 4 ngày làm việc cùng
nhau, công nhân thứ hai nghỉ, công nhân thứ nhất làm thêm 10 ngày nữa mới xong. Hỏi nếu làm riêng
thì mỗi người mất bao nhiêu ngày?
\end{de_cau}

\begin{loi_giai_cau}
Gọi công nhân 1 làm riêng hết $x$ ngày, công nhân 2 hết $y$ ngày ($x,y>0$).

Năng suất: công nhân 1: $\frac{1}{x}$, công nhân 2: $\frac{1}{y}$ (công việc/ngày).

Hệ:
$\frac{1}{x}+\frac{1}{y}=\frac{1}{12}$ và $\frac{4}{x}+\frac{4}{y}+\frac{10}{x}=1$, tức $\frac{14}{x}+\frac{4}{y}=1$.

Đặt $u=\frac{1}{x}, v=\frac{1}{y}$:
\[ \begin{cases}u+v=\frac{1}{12}\\14u+4v=1\end{cases} \].

$14u+4(\frac{1}{12}-u)=1\Rightarrow 10u=\frac{2}{3}\Rightarrow u=\frac{1}{15}\Rightarrow x=15$.

$v=\frac{1}{12}-\frac{1}{15}=\frac{1}{60}\Rightarrow y=60$.

Công nhân 1: \textbf{15 ngày}, công nhân 2: \textbf{60 ngày}.
\end{loi_giai_cau}

\begin{de_cau}{Câu 4}{1,0 điểm}
Cho phương trình $x^2-2(m-1)x+m^2-3m-4=0$. Tìm $m$ để phương trình có hai nghiệm $x_1, x_2$
đều dương và $x_1^2+x_2^2=20$.
\end{de_cau}

\begin{loi_giai_cau}
$\Delta'=(m-1)^2-(m^2-3m-4)=m^2-2m+1-m^2+3m+4=m+5$.

Để có 2 nghiệm dương phân biệt:
$\begin{cases}\Delta'>0\\S=2(m-1)>0\\P=m^2-3m-4>0\end{cases}$

$m>-5$; $m>1$; $(m-4)(m+1)>0\Rightarrow m>4$ hoặc $m<-1$.

Kết hợp: $m>4$.

$x_1^2+x_2^2=S^2-2P=4(m-1)^2-2(m^2-3m-4)=2m^2-2m+12=20$.

$2m^2-2m-8=0\Rightarrow m^2-m-4=0\Rightarrow m=\dfrac{1\pm\sqrt{17}}{2}$.

$m=\dfrac{1+\sqrt{17}}{2}\approx 2{,}56$ (loại vì $<4$). Không có $m$ thỏa mãn.

Kiểm tra lại điều kiện $m>4$: không có nghiệm phù hợp trong tập $m>4$.
\end{loi_giai_cau}
